\section{Related Work}
\label{sec:related_work}

\paragraph{Vector-quantized autoencoders.}
\citet{vandenoord2017vqvae} introduced VQ-VAE by pairing a
convolutional encoder/decoder with a learned codebook and training
the combination with a reconstruction loss and a commitment loss, all
routed through the straight-through estimator.
\citet{razavi2019vqvae2} scaled the architecture to hierarchical
latents and observed that dead codes accumulate without
reinitialisation. VQ-GAN~\citep{esser2021vqgan} combined the VQ
tokenizer with adversarial and perceptual losses and an autoregressive
prior for image synthesis; ViT-VQGAN~\citep{yu2022vitvqgan} improved
both reconstruction quality and codebook usage by replacing the
convolutional backbone with a Vision Transformer and regularising the
codebook lookup. Improved tokenizer training has since been shown to
benefit downstream autoregressive
generation~\citep{sun2024llamagen,yu2024titok,chang2022maskgit} and to
be the critical factor that lets language-model-style generators catch
up with diffusion~\citep{yu2023magvit2}. Our contribution is
orthogonal to these architectural refinements: it modifies only the
gradient flow through the quantizer and leaves the encoder,
decoder, and training losses untouched.

\paragraph{Codebook collapse and rebalancing.}
Codebook collapse, in which a training run converges with the vast
majority of code slots unused, has been attacked from many
angles. Dead-code reinitialisation and EMA updates on the codebook
itself~\citep{vandenoord2017vqvae,razavi2019vqvae2} are standard
remedies; \citet{zheng2023onlineclustered} propose an online-clustering
alternative that re-assigns under-used codes; \citet{baykal2023edvae}
frame the collapse as an overconfidence problem and use evidential
posteriors to spread probability mass across codes; and
\citet{zhang2024optimaltransportvq} enforce a balanced assignment via
optimal-transport matching between encoder outputs and codes. Other
work tackles collapse structurally: \citet{khalil2023llvqvae} replace
a learned codebook with a lattice, \citet{xu2025gaussianvq} build a
Gaussian VAE around the quantizer, and \citet{chen2024vqdimbalance}
analyse the balance between codebook size and embedding dimension.
\citet{takida2024hqvae} train a hierarchical variant whose codebooks
are updated variationally. None of these works change the gradient
estimator itself; our method is therefore complementary and can be
combined with any of them.

\paragraph{Gradient estimation through discrete operators.}
The straight-through estimator of \citet{bengio2013ste}, originally
proposed for stochastic binary neurons, was adapted to VQ-VAEs by
\citet{vandenoord2017vqvae} and remains the default.
\citet{huh2023straighteningste} document the severity of the
STE-induced optimisation pathology for VQ networks and propose
alternating minimisation and a rotation-based residual term that
reshapes the gradient.
\citet{shah2024decoupledste} decouple the forward and backward
identities of the STE to reduce bias.
\citet{liu2023discretebackprop} derive a family of corrections that
interpolate between the STE and the true chain rule. On the
continuous-relaxation side, the
Gumbel-Softmax~\citep{jang2017gumbel} and concrete
distribution~\citep{maddison2017concrete} replace the hard
$\arg\min$ with a temperature-controlled softmax; variants such as
REBAR~\citep[and follow-ups]{jang2017gumbel} trade bias for lower
variance. \citet{vali2022nsvq} substitute random noise for the
quantization residual to obtain an unbiased stochastic estimator.
Our estimator differs in a specific way: its forward output is
bit-exact VQ (no relaxation temperature, no stochasticity) and its
backward Jacobian is the closed-form rotation-plus-scaling that the
forward lookup implicitly performs, rather than the identity, a
correction, or a noise-based surrogate.

\paragraph{Rotation-trick and concurrent work.}
The closest prior art is the \emph{rotation trick} of
\citet{fifty2024rotation}, which independently identified that the
STE discards angular information and proposed to replace the
gradient-path identity with a rotation that maps $\ze$ to $\qvec$.
We differ from~\citet{fifty2024rotation} on three points. First, we
derive the Householder reflection in closed form as
$\rmat = \mathbf{I} - 2\vvec\vvec^\top$ with a specific
unit-norm $\vvec$, yielding an $\mathcal{O}(d)$-per-token
implementation that never materialises the matrix; no additional
parameters are learned. Second, we combine the rotation with a scalar
rescaling $s = \|\qvec\|/\|\ze\|$, and we ablate the two components
separately (\S\ref{sec:results}) to isolate their contributions.
Third, we analyse the numerical-stability regime in which
$\ze \approx \qvec$ and provide a stop-gradient wrapper that
guarantees bit-exact forward output even when the reflection is
ill-conditioned. These changes are orthogonal to the codebook-free
direction taken by FSQ~\citep{mentzer2024fsq}, which eliminates the
lookup entirely by bounding the latent with a tanh and rounding each
channel independently.

\paragraph{Downstream tokenizer-based generation.}
Discrete tokenizers are the backbone of autoregressive and masked
generative transformers that operate in image, video, and
unified-multimodal settings~\citep{esser2021vqgan,chang2022maskgit,
sun2024llamagen,gu2022vqdiffusion,yu2023magvit2,yu2024titok,
kolesnikov2022uvim}. Because these downstream models freeze the
tokenizer during training, any improvement to the tokenizer's
reconstruction fidelity and codebook utilisation propagates directly
into downstream sample quality. Our experiments therefore target
reconstruction metrics first and reserve the downstream generation
benchmark (E5 in our experimental matrix) as a secondary evaluation.
